% Subway Opening in the Process of Urban Revival — LaTeX version of the submitted note (July 2025).
% Text as submitted; the figure is read from the pipeline output, so it updates after `master.py`.
% Build:  cd Writing && latexmk -pdf main.tex
\documentclass[12pt]{article}
\usepackage[a4paper,margin=2.2cm]{geometry}
\IfFileExists{newtxtext.sty}{\usepackage{newtxtext,newtxmath}}{\usepackage{mathptmx}}   % Times
\usepackage[T1]{fontenc}
\usepackage[utf8]{inputenc}
\usepackage{graphicx}
\usepackage{caption}
\usepackage[authoryear,round]{natbib}
\usepackage[hidelinks]{hyperref}
\captionsetup{font=small,labelfont=bf}
\setlength{\parskip}{0.4em}
\setlength{\parindent}{0pt}

\newcommand{\figfile}{../empirics/results/figures/fig1_combined.pdf}

\begin{document}

% ---------------------------------------------------------------- cover page
\begin{titlepage}
\centering
\vspace*{6cm}
{\LARGE Subway Opening in the Process of Urban Revival\par}
\vspace{1.5cm}
{\large Fedor Shafranov\par}
\vspace{0.5cm}
Final assignment, Housing, LMU Munich, July 2025\par
\vfill
{\small Replication package: \texttt{master.py} in the repository root reproduces Figure~1.\par}
\end{titlepage}

% ---------------------------------------------------------------- main page
\section*{Subway Opening in the Process of Urban Revival}

\textbf{Setting:} Hudson Yards in New York City was once underdeveloped compared to other Manhattan
neighborhoods, primarily consisting of unused land and rail yards owned by the MTA. Development began in
2012, becoming one of the largest privately-led real estate projects \citep{fisher2018}. In September 2015,
the 34th Street--Hudson Yards subway station opened. As the last case of a new station built in a completely
new neighborhood in NYC, it offers a unique case study of how commercial development interacts with urban
infrastructure. This study examines how the opening of the subway influences commuter patterns in the area.

\textbf{Data:} To estimate changes in commuting patterns, I use two datasets: the Primary Land Use Tax Lot
Output (PLUTO) from the New York City Department of City Planning \citeyearpar{pluto2025} and the LEHD
Origin-Destination Employment Statistics (LODES) from the U.S. Census Bureau \citeyearpar{lodes2025}.
PLUTO contains data on all tax lots in New York City, updated annually to reflect demolitions and new
constructions. I aggregate the data to the census tract level for the Hudson Yards area to track
neighborhood development. LODES provides a cross-product of census blocks in New York State. I aggregate
this data to the tract level and identify workers commuting to Hudson Yards from each New York census
tract. I then estimate the distance between tracts.

\textbf{Fact:} I observe two key trends in the data. First, the subway opening stimulates construction.
Panel A of Figure~\ref{fig:main} shows that after 2015, development in Hudson Yards accelerated. Second, it
increases commutes for people living farther from the area, replacing those who live nearby. Panel B shows
that, for all distance bins except the closest one, the number of commuters rose.

\textbf{Discussion:} The second fact is of main interest because it shows gentrification of the community
in action. The newly built subway increases construction in the neighborhood and attracts new commuters
because of more available transit and more opportunities.

\textbf{Literature:} The fact relates to previous literature that studies the effects of urban revival and
development of urban infrastructure. Previous literature considered how rail transit development introduced
advanced degree holders to the neighborhood \citep{liang2022}. The other publications analyzed changes in
housing prices \citep{guan2024} and ridership increase \citep{gonzalez2018}. My analysis is novel because
it shows quantitatively how more people living far away from the neighborhood decide to commute because of
new infrastructure and available transit. Previously, changes in commuters were not studied over a
geographic distance.

% ---------------------------------------------------------------- figure
\clearpage
\begin{figure}[p]
\centering
\IfFileExists{\figfile}{\includegraphics[width=\textwidth,height=0.78\textheight,keepaspectratio]{\figfile}}%
  {\fbox{\parbox{0.9\textwidth}{\centering\vspace{3cm}Run \texttt{master.py} to create\\ \texttt{empirics/results/figures/fig1\_combined.pdf}\vspace{3cm}}}}
\caption{Impact of Subway on Construction and Commuting Patterns. Panel A shows changes in construction in
Hudson Yards from 2009 to 2020, with the built area divided into three categories: total area, commercial
space, and office space. The areas are converted from square feet to square meters, and the vertical line
marks the subway opening. Panel B illustrates changes in commuting patterns to Hudson Yards before and after
the subway opened. The y-axis represents the average number of Hudson Yards commuters per 1,000 workers in
their census tract, while the x-axis bins census tracts based on their distance from Hudson Yards.}
\label{fig:main}
\end{figure}

\clearpage
\bibliographystyle{plainnat}
\bibliography{references}

\end{document}
